% Generated draft: requires mathematical and source review.
\documentclass[11pt]{article}
\usepackage[margin=25mm]{geometry}
\usepackage{fontspec}
\setmainfont{Latin Modern Roman}
\usepackage{amsmath,amssymb,mathtools,unicode-math}
\setmathfont{Latin Modern Math}
\usepackage{enumitem,needspace,xurl,hyperref}
\setlength{\parindent}{0pt}
\setlength{\parskip}{5pt}
\setlength{\emergencystretch}{3em}
\newcommand{\N}{\mathbb{N}}
\newcommand{\Z}{\mathbb{Z}}
\newcommand{\Q}{\mathbb{Q}}
\newcommand{\R}{\mathbb{R}}
\newcommand{\C}{\mathbb{C}}
\newcommand{\F}{\mathbb{F}}
\newcommand{\A}{\mathbb{A}}
\newcommand{\PP}{\mathbb{P}}
\newcommand{\E}{\mathbb{E}}
\newcommand{\eps}{\varepsilon}
\newcommand{\dd}{\,\mathrm{d}}
\begin{document}
\section*{UnsolvedMath ID: 30000563}
\textbf{Source problem code:} catalogue entry
\subsection{Definitions and mathematical statement}
Let \(X\) be the smooth closed oriented four-manifold intended in the source, with \(b_2^+(X)>1\). Here \(b_2^+\) is the dimension of a maximal positive subspace of the real intersection form. Fix a homology orientation, conventionally an orientation of \(H^1(X;\R)\oplus H^2_+(X;\R)\), to align signs of the invariants. Let \(f\) be a broken fibration in the intended class. Such terminology usually concerns a map to a surface with indefinite-fold singularities and, in a broken-Lefschetz version, additional Lefschetz singularities; the source does not give its base, allowed singularities, or the hypotheses required for the matching invariant to be defined.

Let \(\mathcal S_f\subseteq\operatorname{Spin}^c(X)\) denote the source's admissible structures, and let \(L_{X,f}(s)\) be the Lagrangian matching invariant. The conjecture is
\[
\forall s\in\mathcal S_f,\qquad
L_{X,f}(s)=SW_X(s),
\]
for every homology-oriented \(X\) with \(b_2^+(X)>1\) and every intended broken fibration \(f\). Both sides must be interpreted in the same coefficient and grading conventions and with corresponding orientation choices.

The record does not define admissibility or specify whether the invariants are numerical counts, graded objects, or require insertions when the expected moduli dimension is nonzero. No restriction to zero-dimensional moduli spaces or reduction modulo two is imposed silently; the homology-orientation hypothesis is retained. The actual definition of \(L_{X,f}\), its auxiliary choices, and the comparison of targets must be recovered.

The stated consequence is independence of the matching invariant from \(f\). For two permitted fibrations \(f,f'\), the equality would give
\[
L_{X,f}(s)=L_{X,f'}(s)
\]
on their common admissible domain, once target identifications are aligned. A claim of a single invariant on all structures additionally requires matching or extending the admissibility domains, which the extract does not specify.

The supplied literature triage checked 2026-08-21 reports the comparison as an unresolved claim and no authoritative resolution found in that check. It supplies no further hypotheses or theorem establishing either equality or full fibration independence.
\subsection{Short English statement}
For the stated four-manifolds and admissible structures, the Lagrangian matching invariant is conjectured to equal Seiberg--Witten theory, giving fibration independence on corresponding domains. Admissibility, target, and orientation conventions remain to be specified.
\subsection{Sources}
\begin{itemize}
\item Local UnsolvedMath registry, permanent ID 30000563, catalogue entry, \emph{Lagrangian Matching and Seiberg--Witten Invariants}.
\item \emph{Four-dimensional Manifolds} (2007), Oberwolfach Reports, pp. 2141--2194, DOI \url{https://doi.org/10.4171/owr/2006/35}.
\item Supplied literature triage checked 2026-08-21 cites only the report; no separate invariant-definition or comparison theorem citation is supplied.
\end{itemize}
\end{document}
